\documentclass[11pt,a4paper]{article}
\usepackage[T1]{fontenc}
\usepackage[utf8]{inputenc}
\usepackage[english]{babel}
\usepackage{lmodern}
\usepackage{microtype}
\usepackage{geometry}
\geometry{margin=2.5cm}
\usepackage{hyperref}
\hypersetup{
    colorlinks=true,
    linkcolor=blue!70!black,
    urlcolor=blue!70!black,
    citecolor=blue!70!black
}
\usepackage{booktabs}
\usepackage{longtable}
\usepackage{enumitem}
\usepackage{xcolor}
\usepackage{titlesec}
\usepackage{fancyhdr}
\usepackage{graphicx}
\usepackage{listings}

\definecolor{codegray}{rgb}{0.95,0.95,0.95}
\definecolor{doctum}{rgb}{0.10,0.30,0.60}

\lstset{
  backgroundcolor=\color{codegray},
  basicstyle=\small\ttfamily,
  breaklines=true,
  frame=single,
  rulecolor=\color{gray!40},
  xleftmargin=1em,
  framexleftmargin=0.5em
}

\pagestyle{fancy}
\fancyhf{}
\fancyhead[L]{\textcolor{doctum}{\small\textbf{Doctum Consilium}}}
\fancyhead[R]{\small\nouppercase{\leftmark}}
\fancyfoot[C]{\thepage}

\titleformat{\section}{\large\bfseries\color{doctum}}{}{0em}{}[\titlerule]
\titleformat{\subsection}{\normalsize\bfseries\color{doctum!80}}{}{0em}{}


\begin{document}

\title{
  \textcolor{doctum}{\Large\textbf{Street Fighter Enhanced (Synology)}} \\[0.5em]
  \large Browser-Based Street Fighter Game — Cocos2d-HTML5 \\[0.3em]
  \normalsize Technical Foundations and Architecture
}
\author{Doctum Consilium \\ \small\textit{Generated 2026-05-08}}
\date{}
\maketitle
\thispagestyle{fancy}

\begin{abstract}
Enhanced browser port of a Street Fighter game built with Cocos2d-HTML5 framework, featuring sprite animation, physics, and Synology NAS deployment.
This document covers the theoretical foundations, architectural decisions,
and key technical concepts required to understand, contribute to, and
maintain this repository.
\end{abstract}

\tableofcontents
\newpage

\section{Overview}

A browser-based 2D fighting game implemented with Cocos2d-HTML5 v3.
The game features sprite-sheet animation, collision detection, physics
simulation, and a multi-layer canvas rendering architecture.
It is containerised with Docker and served as a static web application.


\section{Architecture}

\begin{verbatim}
Browser (Canvas2D / WebGL)
└── Cocos2d-HTML5 runtime
    ├── Scene graph (Layer tree)
    │   ├── BackgroundLayer
    │   ├── AnimationLayer (fighters)
    │   └── StatusLayer (health bars)
    ├── Fighter.js      (state machine)
    ├── FighterPhysics.js (gravity, collision)
    └── Projectile.js   (hadouken physics)
\end{verbatim}


\section{Key Technical Concepts}
\begin{itemize}[leftmargin=1.5em]
  \item Sprite sheet animation: texture atlas (plist) maps named frames to UV coordinates; \texttt{cc.AnimationCache} manages frame sequences.
  \item Scene graph: hierarchical node tree; each node has position, scale, rotation, opacity; parent transforms propagate to children.
  \item Finite state machine (FSM): fighter states (idle, walking, jumping, attacking, stunned, KO); transitions triggered by input or collision events.
  \item AABB collision detection: axis-aligned bounding box overlap; $\text{overlap} = \text{rect}_A.\text{intersects}(\text{rect}_B)$.
  \item Game loop: \texttt{cc.scheduler} drives \texttt{update(dt)} at 60~fps; delta-time normalised physics for frame-rate independence.
\end{itemize}

\section{Data Flow}

Input event (keyboard/touch) →
Fighter FSM transition →
Physics update (gravity, velocity integration) →
Collision detection →
Animation frame selection →
Canvas render.


\section{Technology Stack}

\begin{tabular}{ll}
\toprule
\textbf{Component} & \textbf{Technology} \\
\midrule
Game framework & Cocos2d-HTML5 v3 \\
Rendering & Canvas2D / WebGL \\
Asset format & PNG sprite sheets + plist atlases \\
Server & nginx (Docker) \\
\bottomrule
\end{tabular}


\section{Design Decisions}

\begin{itemize}
  \item Cocos2d-HTML5 chosen for its mature sprite/animation API
    and cross-browser Canvas/WebGL abstraction.
  \item Static file serving via nginx: no backend required,
    minimal operational overhead on Synology NAS.
\end{itemize}


\section{Repository Structure}
\begin{lstlisting}[language=bash,caption={Top-level directory layout}]
streetfighter-enhanced_syno/
  .github/
    copilot-instructions.md   # GitHub Copilot guardrails
    agents/
      ecr-secrets-agent.agent.md  # ECR secret rotation runbook
  CLAUDE.md                   # Claude Code guardrails
  GEMINI.md                   # Gemini CLI guardrails
  INFRASTRUCTURE.md           # Operational runbook
  ROADMAP.md                  # Execution log and roadmap
  README.md                   # Project overview
\end{lstlisting}

\section{Contributing}
\begin{enumerate}
  \item Read \texttt{README.md}, \texttt{INFRASTRUCTURE.md}, and \texttt{ROADMAP.md} before any task.
  \item Follow the Code Documentation Standard: every public function must have
    a docstring (Google-style for Python, JSDoc for TypeScript, XML for C\#).
  \item Run the guardrails validator before committing:
    \begin{lstlisting}[language=bash]
git add -A && ./.guardrails/bin/validate_guardrails.sh
    \end{lstlisting}
  \item For deployment or destructive operations, always use a named tmux session:
    \begin{lstlisting}[language=bash]
tmux new-session -A -s deploy-streetfighter-enhanced_syno
    \end{lstlisting}
\end{enumerate}

\end{document}
